%% Determines the type of document (including standard settings for layout).
\documentclass[letterpaper]{article}



%% Package to control the size of a page: the standard margins used by LaTeX are
%% very wide!
\usepackage[margin=1in]{geometry}


%% Packages add functionality to the document. The AMS packages are standard
%% packages to support various mathematical notations. AMS stands for American
%% Mathematical Society, the organization that maintains these packages.
\usepackage{amsmath,amsthm,amssymb}

%% Theorems-like environments using functionality provided by amsthm.
\theoremstyle{plain}
\newtheorem{theorem}{Theorem}[section]
    %% [section] at the end specifies that theorems should be numbered
    %% per-section: Section x starts with theorem-like x.1 and so on, ...
\newtheorem{proposition}[theorem]{Proposition}
    %% [theorem] in the middle specifies: use the same counter as the theorem
    %% environment: here we number all theorem-like environments consecutively.
\newtheorem{corollary}[theorem]{Corollary}
\newtheorem{lemma}[theorem]{Lemma}
\theoremstyle{definition}
\newtheorem{definition}[theorem]{Definition}
\theoremstyle{remark}
\newtheorem{example}[theorem]{Example}
\newtheorem{remark}[theorem]{Remark}


%% Support for nicely formatted tables.
\usepackage{booktabs}


%% Support for colors & colors in tables.
\usepackage[table]{xcolor}

% Seven colors safe for use color blindness.
% Colors taken from doi:10.1038/nmeth.1618.
\definecolor{cbOrange}{RGB}{230,159,0}
\definecolor{cbSkyBlue}{RGB}{86,180,233}
\definecolor{cbBluishGreen}{RGB}{0,158,115}
\definecolor{cbYellow}{RGB}{240,228,66}
\definecolor{cbBlue}{RGB}{0,114,178}
\definecolor{cbVermillion}{RGB}{213,94,0}
\definecolor{cbReddischPurple}{RGB}{204,121,167}


%% Notation used in this document.
\newcommand{\n}{\mathbf{n}} %% Num. Replicas.
\newcommand{\f}{\mathbf{f}} %% Num. Faulty Replicas.

%% Misc. Math notation.
\newcommand{\BigO}{\mathcal{O}}
\newcommand{\abs}[1]{\lvert #1 \rvert}
\newcommand{\AName}[1]{\textsc{#1}}
\newcommand{\Var}[1]{\texttt{#1}}


%% Algorithms.
\usepackage{algorithm}
\usepackage[noend]{algorithmic}
\newcommand{\GETS}{:=}


%% Formatting SI-units.
\usepackage{siunitx}
\sisetup{per-mode=symbol}


%% TikZ: for creating figures.
\usepackage{tikz}

%% Configuration for figures: Nicer arrows.
\usetikzlibrary{arrows.meta}
\tikzset{>=Stealth}


%% pgfplots: drawing plots using TikZ.
\usepackage{pgfplots}
%% Configuration for plots: Use color-blind friendly colors.
\pgfplotscreateplotcyclelist{cbSafeList}{
    very thick,solid,cbOrange,every mark/.append style={solid},mark=*\\
    very thick,solid,cbSkyBlue,every mark/.append style={solid},mark=*\\
    very thick,solid,cbBluishGreen,every mark/.append style={solid},mark=*\\
    very thick,solid,cbYellow,every mark/.append style={solid},mark=*\\
    very thick,solid,cbBlue,every mark/.append style={solid},mark=*\\
    very thick,solid,cbVermillion,every mark/.append style={solid},mark=*\\
    very thick,solid,cbReddischPurple,every mark/.append style={solid},mark=*\\
    very thick,solid,black,every mark/.append style={solid},mark=*\\
}
\pgfplotsset{
    legend style={font=\small},
    compat=1.16,
    width=260pt,
    height=140pt,
    legend cell align=left,
    xlabel near ticks,
    ylabel near ticks,
    every axis/.append style={
        cycle list name=cbSafeList,
        ymin=0,
        enlargelimits=0.05,
        mark size=1pt,
        ylabel style={align=center},
        xlabel style={align=center},
        title style={align=center}
    }
}

%% PgfplotsTable: loading data files to use with pgfplots.
\usepackage{pgfplotstable}


%% Support for hyperlinks and urls. The setting ``colorlinks'' sets how links
%% are shown in the document (with a color, without underline). We put hyperref
%% last---it has a tendency to break other packages when loaded before them.
\usepackage[colorlinks]{hyperref}
\usepackage{graphicx}
\usepackage{pgffor}
\usepackage{caption}
\usepackage{tabularx}
\usepackage{enumitem}
\usepackage{fancyhdr}


\begin{document}

%% First the meta-data of the doucment: title,

\begin{titlepage}
    \centering
    \vspace*{3cm}
    
    {\Huge COMPSCI 2DB3 Assignment 4}\\[2cm]
    
    {\large Luca Mawyin}\\[0.5cm]
    
    {\large Dr. Jelle Hellings}\\[0.5cm]

    {\large \today}\\[2cm]

    {\Large McMaster University}

    \vfill
    \begin{flushleft}

        {\large \textbf{Student Number: }400531739}\\[0.5cm]
        {\large \textbf{MacID: }mawyinl}
    \end{flushleft}

\end{titlepage}

\stepcounter{section}
\section*{P.\thesection}

$\text{Q}_1:\pi_{uid}(\textbf{Review})~\backslash~\pi_{uid}(\textbf{Recipe})$

\stepcounter{section}
\section*{P.\thesection}

$\text{NOTMOSTRECENT}:=\rho_{C_1.rid\mapsto{rid},C_1.title\mapsto{title},C_1.uid\mapsto{uid},C_1.when\mapsto{when}}(\pi_{C_1.rid,C_1.title,C_1.uid,C_1.when}$\\
$(\sigma_{C_1.uid=C_2.uid\land{C}_1.when<C_2.when}(\rho_{C_1}(\textbf{Recipe})\times\rho_{C_2}(\textbf{Recipe}))))$\\[1.5em]
$Q_2:\textbf{Recipe}~\backslash~\text{NOTMOSTRECENT}$

\stepcounter{section}
\section*{P.\thesection}

$\text{ATLEAST-3}:=\rho_{I_1.rid\mapsto{rid},I_1.ingredient\mapsto{ingredient}}(\pi_{I_1.rid,I_1.\text{ingredient}}\\
(\sigma_{I_1.\text{ingredient}=I_2.\text{ingredient}=I_3.\text{ingredient}\land{I_1.rid}\neq{I_2.rid}\land{I_2.rid}\neq{I_3.rid}\land{I_1.rid}\neq{I_3.rid}}(\rho_{I_1}(\textbf{Ingredient})\times\rho_{I_2}(\textbf{Ingredient})\times\rho_{I_3}(\textbf{Ingredient}))))$\\[1.5em]
$Q3:\pi_{rid,\text{ingredient}}(\textbf{Ingredient})~\backslash~\text{ATLEAST-3}$

\stepcounter{section}
\section*{P.\thesection}

$\text{MUTLI\_MANUFACTURER}:=\rho_{P_1.ingredient\mapsto{ingredient}}(\pi_{P_1.\text{ingredient}}\\
(\sigma_{P_1.\text{ingredient}=P_2.\text{ingredient}\land{P_1.\text{manufacturer}\neq{P_2.\text{manufacturer}}}}(\rho_{P_1}(\textbf{Product})\times\rho_{P_2}(\textbf{Product}))))$\\[1.5em]
$\text{SINGLE\_MANUFACTURER\_ING}:=\pi_\text{ingredient}(\textbf{Product})~\backslash~\text{MULTI\_MANUFACTURER}$\\[1.5em]
$\text{SINGLE\_MANUFACTURER}:=\pi_{rid,\text{ingredient}}((\textbf{Ingredient}\bowtie_{\text{Ingredient.ingredient=SINGLE\_MANUFACTURER\_ING.ingredient}}
\text{SINGLE\_MANUFACTURER\_ING})\bowtie_{\text{Ingredient.ingredient=Product.ingredient}}\textbf{Product})$\\[1.5em]
$Q4:\pi_{rid,\text{manufacturer}}(\textbf{Ingredient}\bowtie_{\text{Ingredient.ingredient=Product.ingredient}}\textbf{Product})~\backslash~\text{SINGLE\_MANUFACTURER}$

\stepcounter{section}
\section*{P.\thesection}

$\text{RA}:=\pi_{rid,author}(\rho_{uid\to\text{author}}(\textbf{Recipe}))$\\
$\text{RV}:=\pi_{reviewer,rid}(\rho_{uid\to\text{reviewer}}(\textbf{Review}))$\\
$\text{RECIPEREVIEWS}:=\pi_{\text{author,reviewer},rid}(\text{RV}\bowtie_{RV.rid=RA.rid}\text{RA})$\\[1.5em]
$\text{Q5}:\text{RECIPEREVIEWS}\div\text{RA}$

\stepcounter{section}
\section*{P.\thesection}

$\text{RI}:=\pi_{rid,\text{ingredient}}(\textbf{Ingredient})$\\
$\text{SI}:=\rho_{sid\mapsto{rid}}(\pi_{sid,\text{ingredient}}(\text{RI}))$\\
$\text{TI}:=\rho_{tid\mapsto{rid}}(\pi_{tid,\text{ingredient}}(\text{RI}))$\\
$\text{MATCHINGREDIENTS}:=(\text{SI}\bowtie_{\text{SI.ingredient}=\text{TI.ingredient}\land\text{SI}.sid\neq\text{TI}.tid}\text{TI})$\\[1.5em]
$\text{Q6}:\pi_{sid,tid}(\text{MATCHINGREDIENTS}\div\pi_{\text{ingredient}}(\text{SI}))$
\end{document}